\subsection{Raw Bitmap Representation}

A raw RGB bitmap stores the colour components of the pixels directly.
For a $2\times2$ group of pixels,
\begin{align}
\begin{bmatrix}
(R_{11},G_{11},B_{11}) & (R_{12},G_{12},B_{12})\\
(R_{21},G_{21},B_{21}) & (R_{22},G_{22},B_{22})
\end{bmatrix}\label{eq:raw_rgb_2x2}
\end{align}

\vspace*{\baselineskip}

\subsection{JPEG Representation}

JPEG uses a transformed and compressed representation rather than
storing every RGB tuple independently. The JPEG standard defines processes
for converting source image data to compressed image data.

The common colour information is calculated, and the formulae used for
this conversion are listed in the next section.

\vspace*{\baselineskip}

\subsection{Per-Pixel Storage}

For an image containing $N$ pixels, an 8-bit RGB representation requires
three separate 8-bit colour components for every pixel. Hence, the raw
colour information requires
\begin{align}
3\times8N&=24N\label{eq:rgb_storage_bits}
\end{align}
bits.

If the image is converted directly from RGB to grayscale, only the
luminance value $Y$ is retained for each pixel. With an 8-bit grayscale
representation, this requires
\begin{align}
1\times8N&=8N\label{eq:grayscale_storage_bits}
\end{align}
bits.

Thus, the storage changes from $24N$ bits for RGB to $8N$ bits for
grayscale. The storage ratio is therefore
\begin{align}
\frac{8N}{24N}&=\frac{1}{3}\label{eq:rgb_grayscale_ratio}
\end{align}
so the grayscale representation requires one-third of the raw RGB
storage, corresponding to a reduction of two-thirds.

When the image is represented using $Y$, $Cb$, and $Cr$, the three
components are still separate components. Each component is normally
represented using 8-bit samples; the reduction in storage can instead be
obtained by reducing the spatial sampling of the chrominance components.
In particular, with 4:2:0 chroma subsampling, the $Y$ component retains
$N$ samples, while $Cb$ and $Cr$ each contain $N/4$ samples.

Therefore, before the additional JPEG transform and coding stages, the
sample storage can be written as
\begin{align}
8N+8\frac{N}{4}+8\frac{N}{4}
&=12N\label{eq:ycbcr_420_storage_bits}
\end{align}
bits.

Thus, the factor of four applies to the \emph{number of chrominance
samples}, not to the number of bits in an individual $Cb$ or $Cr$ sample.
The $Y$ component is retained at full spatial resolution because it carries
the intensity information for every pixel.

For example, in a $2\times2$ group containing four $Y$ samples, 4:2:0
subsampling uses one $Cb$ sample and one $Cr$ sample for the group. The
$Cb$ and $Cr$ samples are therefore stored separately, but they are shared
spatially by the four corresponding pixels.

The comparison above describes sample storage before the main JPEG
compression stages. The actual size of a JPEG file is different because
JPEG additionally uses transformation, quantisation, and entropy coding
as part of its coding process.
